% % % % % % % % % % % % % % % % % % % % % % % % % % % % % % % % % % % % % % % % % % % % % % % % % % % % % % % % % % % % % % % %
% SAT_STYLE_MATH_MCQS.tex
% 1_09.tex
% 
% encoding: UTF-8 
% EOL: LF
%
% format: LaTeX
% indent: spaces (2)
% width: 127
% % % % % % % % % % % % % % % % % % % % % % % % % % % % % % % % % % % % % % % % % % % % % % % % % % % % % % % % % % % % % % % %
Considérons trois nombres $a, b, c$ dont la moyenne vaut $150$. Sachant que la moyenne de $a$ et $b$ vaut $120$ et que %
$a - b = 40$, que vaut le triplet $(a, b, c)$?
%
\begin{enumerate}
  \item $(130,90,70)$
  \item $(60,100,350)$
  \item $(140,100,210)$
  \item $(140,100,70)$
  \item $(120,80,210)$
\end{enumerate}
%
\begin{proof}
Analogue à l'[\ref{1.8}], ce problème se formalise aussi commme un système de trois équations à trois inconnues $a, b, c$, %
comme suit:
%
\begin{equation}
  \left.
  \begin{aligned}
    (a+b+c)/3 &= 150\,\, \\
    (a+b)/2   &= 120 \\
    a - b           &= 40 \\
  \end{aligned}
  \right\} \iff
  \begin{cases} \label{1.9:equations}
    a+b+c &= 450 \\
    a+b   &= 240 \\
    a - b &= 40 
  \end{cases}
\end{equation}
%
Ajouter la deuxième ligne à la troisième donne $a$. On peut alors trouver $b$:
%
\begin{align}
  \left.
  \begin{aligned}
    a+b   &= 240\,\, \\
    a - b &= 40 
  \end{aligned}
  \right\} & \iff 
  \begin{cases}
    a = 140 \\
    b = 100. \\
  \end{cases}
\end{align}
%
Revenant à la première ligne, nous concluons:
%
\begin{align}
  \left.
  \begin{aligned}
    a + b + c &= 450\,\, \\
    a &= 140 \\
    b &= 100
  \end{aligned}
  \right\} & \iff
  \left\{
  \begin{aligned}
    \hspace{0.25em}\boxed{c = 210} \\
    \hspace{0.25em}\boxed{a = 140} \\
    \hspace{0.25em}\boxed{b = 100}
  \end{aligned}
  \right.
\end{align}
%
Cette résolution, par le \textit{pivot de Gauss}, peut être ``codée'' sous la forme d'une matrice $M$, comme suit:
%
\begin{align*}
  \begin{bNiceArray}{cccc}[
    %margin, 
    extra-margin=11.5pt, 
    cell-space-limits=5pt
  ]
    \mathtt{1} & \mathtt{1} & \mathtt{1} & \mathtt{450} \\
    \mathtt{1} & \mathtt{1} &  & \mathtt{240} \\
    \mathtt{1} & \texttt{-1} &  & \hspace{0.5em}\mathtt{40}
  \end{bNiceArray}  & &  \\
  \begin{bNiceArray}{cccc}[
    %margin, 
    extra-margin=14pt, 
    cell-space-limits=5pt
  ]
  \mathtt{1} & \mathtt{1} & \mathtt{1} & \mathtt{450} \\
  \mathtt{1} & \mathtt{1} &  & \mathtt{240} \\
  \mathtt{2} &            &  & \mathtt{280}
  \end{bNiceArray}  & \hspace{1cm}\texttt{M[2] = M[1] + M[2]} & \texttt{\% M[i] est la ligne i}\\
%
  \begin{bNiceArray}{cccc}[
    %margin, 
    extra-margin=8pt,
    cell-space-limits=5pt
  ]
  \mathtt{1} & \mathtt{1} & \mathtt{1} & \mathtt{450} \\
  \mathtt{1} & \mathtt{1} &            & \mathtt{240} \\
  \mathtt{\boxed{1}} &            &            & \mathtt{\boxed{140}}
  \end{bNiceArray} & \hspace{1cm}\texttt{M[2] = M[2] / 2} & \texttt{\% a = 140} \\
%
  \begin{bNiceArray}{cccc}[
    %margin, 
    extra-margin=8pt,
    cell-space-limits=5pt
  ]
  \mathtt{1} & \mathtt{1} & \mathtt{1} & \mathtt{450} \\
             & \mathtt{\boxed{1}} &            &  \mathtt{\boxed{100}} \\
  \mathtt{1} &            &            &  \mathtt{140}
  \end{bNiceArray}  & \hspace{1cm}\texttt{M[1] = M[1] - M[2]} & \texttt{\% b = 100} \\
%
  \begin{bNiceArray}{cccc}[
    %margin, 
    extra-margin=8pt,
    cell-space-limits=5pt
  ]
             &             & \mathtt{\boxed{1}} &  \mathtt{\boxed{210}} \\
             & \mathtt{1}  &                    &  \mathtt{100} \\
  \mathtt{1} &             &                    &  \mathtt{140}
  \end{bNiceArray}  & \hspace{1cm}\texttt{M[0] = M[0] - M[1] - M[2]} & \texttt{\% c = 210}
\end{align*}
Réponse \textbf{C}.
%
\end{proof}
%